
\def\PUDcodigo{ACE004}

\if\COMPILAR0
    \def\PUDcodacad{04507.3}
\else
    \def\PUDcodacad{04506.3}
\fi

\def\PUDdisciplina{Inglês Técnico}

\def\PUDsemestre{S1}

\def\PUDhoras{40}

%NOVOS ---------------------------------------------------
\def\PUDhorasTeoria{40}
\def\PUDhorasPratica{0}

\def\PUDinterECA{---}

\def\PUDatuaisECA{- Utilização de textos diversos, com ênfase em textos técnicos contendo temas atuais da área industrial.}

\def\PUDrecursos{Projetor multimídia; Quadro branco e pincel.}

%----------------------------------------------------------

\def\PUDcreditos{4}

\def\PUDprerequisito{---}

\def\PUDpreacad{---}

\def\PUDobjetivos{Enriquecer o vocabulário em língua inglesa; Aprimorar a capacidade de compreensão de textos diversos, com ênfase em textos técnicos afins à atividade profissional.}

\def\PUDementa{Desenvolvimento da habilidade de retirar informações fidedignas e relevantes de textos técnico-científicos autênticos, redigidos em língua inglesa;  Conscientização das estratégias de processamento textual superficiais e profundas, visando ao desenvolvimento da habilidade de leitura;  Consolidação das estruturas gramaticais típicas do discurso acadêmico.}

\def\PUDconteudo{ &  Considerações gerais sobre leitura; 
\\ 
 &  Estrutura da frase em Língua Inglesa; 
\\ 
 &  Introdução às estratégias de leitura; 
\\ 
 &  Lay-out;  \\ 
 &  Skimming/scanning; 
\\ 
 &  Utilização de informação não-linear; 
\\ 
 &  Key words; 
\\ 
 &  Congnates; 
\\ 
 &  Word formation; 
\\ 
 &  Linking Word; 
\\ 
 &  Interpretação dos marcadores de discurso. }

\def\PUDmetodo{Aulas expositórias que podem ser teóricas e/ou práticas, onde as práticas no laboratório serão marcadas no decorrer da disciplina.}

\def\PUDavalia{A avaliação será feita com aplicação de provas teóricas e/ou práticas, além da possibilidade de inclusão de trabalhos e seminários no decorrer da disciplina.}

\def\PUDbibbasica{ &  \href{www.biblioteca.ifce.edu.br/index.asp?codigo_sophia=24073}{Almeida}, Rubens Queiroz de.  As palavras mais comuns da Língua Inglesa: Desenvolva sua Habilidade de Ler Textos em Inglês.  São Paulo, SP: Novatec, 2002.  ISBN: 8575220373.
\\ 
 &  \href{www.biblioteca.ifce.edu.br/index.asp?codigo_sophia=24413}{Hornby}, A. S.  Oxford Advanced Learner’s Dictionary of Current English.  7º ed. Oxford: Oxford University Press, 2007.  ISBN: 9780194001168.
\\ 
 &  \href{www.biblioteca.ifce.edu.br/index.asp?codigo_sophia=23802}{Almeida}, Rubens Queiroz de.  Read in English: Uma Maneira Divertida de Aprender Inglês.  São Paulo, SP: Novatec, 2002.  ISBN: 8575220225. }

\def\PUDbibcomplementar{ &  \href{www.biblioteca.ifce.edu.br/index.asp?codigo_sophia=23925}{Murphy}, Raymond.  Essential grammar in use: a self-study reference and practice book for elementary students of English: with answers.  3º ed.  New York, USA: Cambridge University Press, 2007.  ISBN: 9780521675819.
\\ 
 &  \href{www.biblioteca.ifce.edu.br/index.asp?codigo_sophia=23826}{Munhoz}, Rosângela.  Inglês Instrumental: estratégias de leitura – módulo II.  São Paulo, SP : Textonovo, 2004.  136p.  ISBN: 858573440-X.
\\ 
 &  \href{www.biblioteca.ifce.edu.br/index.asp?codigo_sophia=24251}{Lopes}, Carolina.  Inglês instrumental: leitura e compreensão de textos.  Fortaleza, IFCE, 2012.  ISBN: 9788564778016.
\\
 &  \href{www.biblioteca.ifce.edu.br/index.asp?codigo_sophia=40056}{Glendinning}, Eric H.  Basic english for computing.  Oxford (Inglaterra): Oxford University Press, 2012.  136p.  ISBN: 9780194574709.
\\
 &  \href{www.biblioteca.ifce.edu.br/index.asp?codigo_sophia=40048}{Glendinning}, Eric H.  Oxford english for information technology.  2º ed.  New York, USA: Oxford University Press, 1998.  222p.  ISBN: 9780194574921. }

\def\PUDautor{Samuel Vieira Dias}

\def\PUDdata{2013-04-23}

\def\PUDrevisao{1}

\def\PUDrevisor{Samuel Vieira Dias}

\def\PUDdatarevisao{2017–05-09}

\PUDcorpo
